\section{Proofs for Discrete-Logarithm Based \AVSS}
\label{appendix:avss-proofs}

In Section~\ref{avss:concrete},
we define a concrete \AVSS which builds upon \FeldVariant and is secure in the honest majority setting.
We demonstrate the correctness,
aggregated secrecy,
and uniqueness of that scheme now.

\subsection{Correctness}

The honest majority \AVSS scheme of Section~\ref{avss:concrete} fulfills the notion of correctness for the base VSS scheme,
as the algorithms $\issueshares, \recover, \verify, \derivepshare$ are identical to \FeldVariant,
as is the map $\vfunc$.
We now demonstrate that the concrete \AVSS scheme fulfills the additional notion of correctness for an \AVSS scheme.

\begin{theorem}
  The honest majority \AVSS scheme of Section~\ref{avss:concrete} is correct in the sense of Figure~\ref{fg:avss-correctness}.
\end{theorem}


\begin{proof}

We analyze each case where a return value of 0 could occur in Figure~\ref{fg:avss-correctness}, and demonstrate that that event will not occur.

\medskip

\begin{case} \label{case:corr-verify}
We first show that for each $i \in \set{n}$, $\avss.\verify(i, \ashare_i, \avsscommit) = 1$.
Recall that $\ashare_i$ is simply
  \[ \ashare_i = \tweak + \sum_{j \in \set{\total}} \share_{ji} \]
by Equation~\ref{eq:agg-share-1}.
Because each $\share_{ji}$ was generated by the $\avss.\issueshares$ algorithm, then $\share_{ji} = f_j(i)$ for some polynomials $f_1,\dots,f_\total$ each of degree $t-1$.
Let $f' = \tweak + \sum_{i \in \set{\total}} f_i$.
Then $(i,\ashare_i)$ is a valid point on $f'$, because
  \[ \ashare_i = \tweak + \sum_{j \in \set{\total}} \share_{ji} =  \tweak + \sum_{j \in \set{\total}} f_j(i) = f'(i) \]
Because $\avsscommit$ is a commitment to the same $f'$ as shown in Equations~\ref{eq:const-agg} and \ref{eq:non-const-agg},
then $\avss.\verify(i, \ashare_i, \avsscommit)$ will output $1$.
Note also that $\avss.\recover$ will output $\aggsk = f'(0) = \tweak + \sum_{j \in \set{\total}} s_j$.
\end{case}

\medskip

\begin{case} \label{case:sk-verify}
We next prove that $\avss.\verify(0, \aggsk, \avsscommit) = 1$.
As above,
after performing $\avss.\aggregatepub$,
the resulting aggregated commitment $\avsscommit$ can be represented as a commitment to the polynomial $f'$.
$\avss.\verify(0, \aggsk, \avsscommit)$ checks that $g^{\aggsk} = \hat{A}_0$ (from Eq.~\ref{eq:vss-verify} with $i=0$).
From Eq.~\ref{eq:const-agg}, $\hat{A}_0 = g^{\tweak} \cdot \prod_{j \in \set{\total}} \vsscommit_j[0] = g^{\tweak} \cdot \prod_{j \in \set{\total}} g^{f_j(0)} =  g^{\tweak} \cdot \prod_{j \in \set{\total}} g^{s_j}$.
Since $\aggsk = \tweak + \sum_{j \in \set{\total}} s_j$ as above,
$\avss.\verify(0, \aggsk, \avsscommit)$ will output 1.

%Recall that $a_{ij}, j \in \set{t-1}$ are generated when performing $\avss.\issueshares$,
%and serve as the coefficients for the polynomial $f_i$ whose constant term is $s_i$,
%as shown in Section~\ref{verifiable-ss}.
%The resulting aggregated secret $\aggsk = f'(0)$ is defined with respect to $\vfunc$, which for the concrete scheme, is addition.
%
%\begin{align} \label{eq:joint-polynomial-commit}
  %\begin{split}
    %& g^{f'(x)} = \prod_{k=0}^{t-1} \avsscommit[k]^{x^k} = g^{(\tweak + \sum_{j=1}^\total a_{j0}) x^0} \cdot g^{\sum_{j=1}^\total a_{j1}x^1} \cdot  \ldots \cdot g^{\sum_{j=1}^\total a_{j(t-1)}x^{t-1}}, \text{ or} \\
    %& \aggsk = f'(0) = \vfunc(f_1(x), \ldots, f_\total(x), \tweak) = \tweak + f_1(x) + \ldots + f_\total(x)
  %\end{split}
%\end{align}
%
%Hence, $\avss.\verify(0, \aggsk, \avsscommit)$ will output 1.
\end{case}

%By extension,
%Case~\ref{case:sk-verify-vfunc} also proves that $\aggsk \neq \avfunc(\{ s_j \}_{j \in \total})$.
%Recall the definition of $\avfunc$ from Equation~\ref{eq:avfunc-conc}.

%We finally demonstrate that
%$\avss.\recover(t, \{(k,\ashare_k)\}, k \in C, \avsscommit) \rightarrow \aggsk$.
%
%\begin{case} \label{case:corr-secret}
%Let $\recoveryset = \{(k,\ashare_k)\}, k \in C$ such that $C \subseteq \set{n}$, $\lvert C \rvert \geq t$.
%
%As in Case~\ref{case:corr-verify}, each of the at least $t$ points in $\recoveryset$ lies on the degree $t-1$ polynomial $f'$, and so $\recover$ will interpolate $f'(0) = \aggsk$.
%%Again, let $f'(0) = \tweak + \sum_{j \in \set{\total}} s_j$, as shown ``in the exponent'' in Equation~\ref{eq:const-agg}.
%%Because each $\ashare_i$ is valid with respect to $\avsscommit$ as shown in Case~\ref{case:corr-verify},
%%then $\ashare_i$ is therefore the point $f'(i)$ on the $t-1$ degree polynomial $f'$.
%%Hence, combining $t$ distinct aggregated shares (in other words,  $t$ distinct points) will uniquely define $f'$,
%%by polynomial interpolation.
%%And so because $\avss.\recover$ performs polynomial interpolation to evaluate $f'(0)$,
%%when provided with a valid $\recoveryset$ as its input.
%%$\avss.\recover$  will output $f'(0) = \aggsk$.
%\end{case}
%
\noindent This concludes the proof.
\end{proof}

\begin{comment}
\subsection{Aggregated Security}

Because in the concrete scheme,
there is a one-to-one correspondence between secret and public values,
then aggregated secrecy implies aggregated security (see Remark~\ref{rem:agg-sec-implies-security}).
Hence, we simply prove aggregated secrecy, next.

\end{comment}

\subsection{Aggregated Secrecy}

We now show that the concrete scheme of Section~\ref{avss:concrete} fulfills the notion of aggregated secrecy from Figure~\ref{fg:avss-secrecy}.
To do so, we first show how the environment simulates $\simshare$ and $\Hash$.
Without loss of generality, we assume $\lvert \corrupt \rvert = t-1$.

\begin{figure} [!ht]
  \begin{pchstack}[boxed,center,space=1em]
  \procedure[linenumbering]{$\simshare(\lambda, \corrupt, \vsschal)$}{
    \secretone \rgets \Zq;\ A_0 = \vsschal \cdot g^{\secretone} \\
    \pcfor j \in \corrupt \pcdo \\
    \pcind \share_j \rgets \Zq \\
    \pcfor k \in \{1, \ldots, t-1 \} \pcdo \\
    \pcind A_k = A_0^{L_{0,k}'} \cdot \prod_{j \in \corrupt} g^{L_{j,k}' \cdot \share_j} \\
   \vsscommit = \langle A_0, \ldots, A_{(t-1)} \rangle \\
   \pcreturn (\secretone, \{(j,\share_j)\}_{j \in \corrupt}, \vsscommit)
}
\pchspace
  \procedure[linenumbering]{$\Hash(\pubaggregateset_i, \auxstring_i)$}{
    \text{\algcom $Q_1,Q_2,Q_3$ are defined} \\
    \text{\algcom by the experiment} \\
    \pcif  Q_3[\pubaggregateset_i, \auxstring_i] \neq \bot \\
    \pcind \pcreturn Q_3[\pubaggregateset_i, \auxstring_i] \\
    \vsscommit^* \gets Q_1 \\
    \{ \vsscommit_j \}_{j \in \set{\simtotal}} \gets Q_2 \\
    \pubaggregateset = \{ \vsscommit_j \}_{j \in \set{\simtotal}} \cup \vsscommit^* \\
    \pcif \auxstring_i \neq \auxstring \textbf{ or } \pubaggregateset_i \neq \pubaggregateset \\
    \text{\algcom Respond honestly if this is} \\
    \text{\algcom not the challenge query} \\
    \pcind  Q_3[\pubaggregateset_i, \auxstring_i] \rgets \Zq \\
    \pcind \pcreturn Q_3[\pubaggregateset_i, \auxstring_i] \\
    \secretone \gets Q_1 \\
%
    \pcfor \{ (j, \share_{kj}) \}_{j \in \honest, k \in \set{\simtotal}} \in Q_2 \\
    \pcind \mathsf{in} \gets (t, \{ (j, \share_{kj}) \}_{j \in \honest}) \\
    \pcind s_k \gets \avss[\Hash].\recover(\mathsf{in}) \\
%
    \tweak \gets (-1) \cdot (\secretone + \sum_{i \in \set{j}} s_k) \\
%
   Q_3[\pubaggregateset_i, \auxstring_i] \gets \tweak \\
   \pcreturn \tweak
}
  \end{pchstack}
  \caption{Algorithms to simulate secret sharing and tweak generation in an \AVSS to an adversary.
  %controlling $t-1$ participants denoted by the set $\corrupt$ with respect to a challenge $\vsschal$ and at least $t$ honest players denoted by the set $\honest$, where $\corrupt \cup \honest = \set{n}$ and $\corrupt \cap \honest = \emptyset$.
  %Here, $L_{i,k}'$ denotes the $k\supth$ coefficient of the Lagrange polynomial $L_i'(x)$,
  %which is computed with respect to the set $\{0\} \cup \corrupt$.
  }
  \label{fg:simulate-avss}
\end{figure}



\begin{lemma} \label{lem:avss-simshare}
  When playing the secure aggregation game as in Figure~\ref{fg:avss-secrecy} against the concrete scheme,
  the output of $\getshare$ when $b=0$ (i.e., when the environment honestly performs $\avss.\issueshares$) is indistinguishable to an adversary from the output when $b=1$ (i.e., when the environment performs $\simshare$ as shown in Figure~\ref{fg:simulate-avss}),
  assuming the adversary controls no more than $\lvert \corrupt \rvert = t -1$ participants.
\end{lemma}
%\ian{What exactly do you mean by this? $\simshare$ has different inputs and outputs from $\issueshares$ ($\secretone$ is an output of $\simshare$ and the secret is an input to $\issueshares$), so what concretely do you mean by ``indistinguishable''?}
% CK- re-worded. This is in the context when the adversary queries the getshare oracle.


\begin{proof}
  $\simshare$ accepts a challenge $\vsschal \in \GG$,
  and outputs a randomly sampled secret $\secretone$,
  $\lvert \corrupt \rvert$ shares,
  and an \AVSS commitment.
  To begin, it picks $\secretone \in \Zq$ at random,
  and then generates a blinded commitment $A_0$ to the (unknown) discrete logarithm of $\vsschal$ using $\secretone$ as the blinding factor, by Equation~\ref{eq:const-coeff-comm}.

  \begin{equation} \label{eq:const-coeff-comm}
    A_0 = \vsschal \cdot g^{\secretone}
  \end{equation}

  $\simshare$ then picks $t-1$ shares for the corrupted parties at random,
  resulting in the list of shares $\langle\share_j\rangle_{j \in \corrupt} \rgets \Zq^{(t-1)}$ and then generates commitments to coefficients
  $A_1, \ldots, A_{(t-1)}$ with respect to these shares and $A_0$, by Equation~\ref{eq:derive-coeff}.

  \begin{equation} \label{eq:derive-coeff}
    A_k = A_0^{L'_{0,k}} \cdot \prod_{j \in \corrupt} g^{L'_{j,k} \cdot \share_j}
  \end{equation}

  %\ian{This math doesn't look right to me, and the proof below doesn't use this equation at all, but just handwaves about how this is committing to $f$ in the exponent.  Please check this math (and everywhere the similar ``coefficients of Lagrange polynomials'' technique is used).}
  %\chelsea{See the extra information I added to the prelim under polynomial interpolation. the math is correct.}

  Here, $L_{i,k}'$ denotes the $k\supth$ coefficient of the Lagrange polynomial $L_i'(x)$,
  which is computed with respect to the set $\{0\} \cup \corrupt$.
  The commitment $\vsscommit$ to the $\lvert \corrupt \rvert$ shares is then $\vsscommit = \langle A_0, \ldots, A_{(t-1)} \rangle$.
  The simulation is perfect because $\vsscommit$ simply is committing to a random polynomial $f$ ``in the exponent'',
  whose constant term is the discrete logarithm of $A_0$,
  (where  $A_0$ is a blinded commitment to the discrete logarithm of $\vsschal$)
  and with $t-1$ random points that are the shares distributed to the adversary.
  Each additional coefficient $a_k,  k \in \set{t-1}$ of $f$ is committed to ``in the exponent'' via $A_k$,
  which can be determined via Equation~\ref{eq:derive-coeff}.
  Hence, so long as the adversary controls fewer than $t-1$ players,
  $\avss.\verify(i, \share_i, \vsscommit)$ for all $i \in \corrupt$ will output $1$.
  \qed
\end{proof}

\begin{lemma} \label{lem:avss-programtweak}
  The adversary has negligible advantage in distinguishing the environment's simulation of $\gettweak$ in the concrete scheme of Section~\ref{avss:concrete} in the (programmable) random oracle model, assuming an honest majority.
\end{lemma}


\begin{proof}
  Recall that $\gettweak$ is instantiated by $\Hash$ in the concrete scheme.
  For any query $(\pubaggregateset', \auxstring' )$ where $\auxstring_i \neq \auxstring$,
  the environment simply programs $\Hash$ to return a random value.
  However, when it receives a query for $(\pubaggregateset, \auxstring)$,
  it programs $\Hash$ to return $\tweak$ as in Equation~\ref{eq:programmed-tweak}.
  Recall that $\auxstring$ is provided to the adversary only after it has finished querying $\getshare$ and $\recvshare$,
  and so similarly $\pubaggregateset$ is fixed at the time that the adversary can query with $\auxstring$ as input
  (i.e, the environment will never have to guess which query to program).

  \begin{equation} \label{eq:programmed-tweak}
   \tweak \gets -\secretone - \sum_{j \in \set{\simtotal}} \secret_j \\
  \end{equation}
  %\ian{This isn't quite what Figure 6 is computing; the first element of $Q_1$ is $\secretone$, but that's already one of the $\total$ $\secret_j$ values, right? Also as commented in Figure 5, $\total$ is actually the sum of the sizes of $Q_1$ and $Q_2$, which is correct for Eq.~\ref{eq:programmed-tweak}, but incorrect in Figure 6.}
  %\chelsea{$\secretone$ is not one of the $\total$ $\secret_j$ values; it is handled seperately}
  %\chelsea{the variable $\total$ was being overloaded. Now, $\total$ means the number of shares generated and received. $\simtotal$ is $\total-1$, as it excludes counting the challenge value.}

  There are three requirements that must be satisfied for the environment to program $\Hash$ correctly.
  \begin{enumerate}
    \item The programmed output $\tweak$ must be indistinguishable from a random value.
    \item The environment must be able to derive $\{ \secret_j \}_{ j \in \set{\simtotal}}$ without the involvement or detection of $\adv$.
    \item The environment can correctly program $\Hash$ strictly before $\adv$ can query on inputs $(\pubaggregateset, \auxstring)$.
  \end{enumerate}

  We first demonstrate that the output $\tweak$ from Equation~\ref{eq:programmed-tweak} is indistinguishable from any other output from $\Hash$.
  Recall that $\secretone$ is chosen at random when $\simshare$ is performed,
  and published only within the commitment $A_0$ as shown in Equation~\ref{eq:const-coeff-comm},
  serving as a blinding factor for $\vsschal$.
  Because $\vsschal$ is chosen at uniformly at random from the public domain,
  and because $\secretone$ is chosen at random,
  the $\tweak$ value in Equation~\ref{eq:programmed-tweak} is indistinguishable from a random element to $\adv$.
  %serves essentially as a ``one-time-pad,''
  %as $\secretone$ is used exactly once during the experiment,
  %\ian{But that's not true: it's used twice: once to create $A_0$, and once to create $\tweak$. And here's where the question about the unbound variable $\vsschal$ in Def~6 is critial! If the adversary knows something about the distribution of $\vsschal$ (and it's not uniform), it \emph{can} distinguish $b=0$ from $b=1$, since in the $b=1$ case, $\Htweak$ will output a tweak that will make the resulting public key (which the adversary will learn) equal to $\vsschal$.}
  % CK- changed wording above, please edit if you think this explanation is insufficient.
%
  Moreover, even if the adversary submitted inputs to $\recvshare$ in the attempt to cancel out the outputs from $\getshare$, the adversary would be unable to do so,
  because the adversary sees only the blinded commitment $A_0$ as shown in Equation~\ref{eq:const-coeff-comm},
  and because $\vsschal$ is chosen uniformly at random.
  Therefore, programming $\Hash$ at exactly one point with $\tweak$ is indistinguishable to any other point that is honestly programmed with a random value chosen from $\Zq$.

  Second, the environment can derive $\{ \secret_j \}_{ j \in \set{\simtotal}}$ without the involvement or detection of $\adv$,
  because we operate in the honest majority setting, and assume that the environment simulates at least $t$ honest players. The environment can therefore reconstruct the adversary's $\secret_j$ values (from $Q_2$), and it knows its own (from both $Q_1$ and $Q_2$).
%  \ian{That's not what's in $Q_1$ any more. Doesn't it get its own $\secret_j$ values from $Q_2$ as well now? Wait: $Q_2$ is weird: in $\getshare$, \emph{all} shares are added to $Q_2$, not just the honest ones, but in $\recvshare$, it's just the honest ones.  Should it be just the honest ones in $\getshare$ as well? Do you never need the corrupt shares later?}
  % CK- it doesn't matter if the environment gets additional shares as recover will output the correct value either way. But I'll change it so it is less confusing.

  Finally, the adversary does not learn $\auxstring$ until \emph{after} it has completed its queries to $\getshare$ and $\recvshare$.
  Before returning $\auxstring$ to $\adv$,
  the environment programs $\Hash$ with $\tweak$ on inputs $(\pubaggregateset, \auxstring)$.
  Because $\auxdistr = \Zq$ and $\auxstring$ is some uniform element from $\auxdistr$,
  then $\adv$ had negligible probability in learning the output $\Hash$ with input
  $\auxstring$ until the environment has programmed it.
  The environment is able to correctly program on these inputs using $\tweak$ because it is guaranteed that $\adv$ cannot further influence $\pubaggregateset$.
  In other words, there is no danger that the environment must make a decision on which input to program using $\tweak$.
  %\ian{This is a good argument, but this is exactly where you're relying on the structure of $\auxdistr$, in particular that it has cryptographic guessing entropy.  If $\auxdistr$ had low entropy, than $\adv$ \emph{could} guess it and issue calls to $\Htweak$ before $\isterminated=1$, breaking the scheme. So you need to say what $\auxdistr$ is in the definition of the concrete scheme, or at the very least, say that it can be anything with at least $\lambda$ bits of guessing entropy.}
  % CK- added this specification in the definition and that observation here.

  This completes the proof.
  \qed
\end{proof}

\begin{theorem} \label{thm:avss-secrecy}
  The concrete scheme of Section~\ref{avss:concrete} fulfills aggregated secrecy against an adversary $\adv$ as in Figure~\ref{fg:avss-secrecy}, assuming an honest majority.
\end{theorem}


\begin{proof}
  Recall that that in the aggregated secrecy game,
  $\adv$ has three opportunities to distinguish the $b=0$ game from the $b=1$ game:
  \begin{enumerate}
    \item by how the environment responds to $\getshare$,
    \item by how the environment responds to $\Ogettweak$, and
    \item whether or not
  $\avss.\derivepshare(0, \avsscommit)$ evaluates to the challenge $\vsschal$ when $b=1$.
  \end{enumerate}

  For $(1)$, by Lemma~\ref{lem:avss-simshare},
  $\getshare$ is indistinguishable to the adversary when $b=0$ (i.e., shares are honestly generated)
  and when $b=1$ (i.e., $\simshare$ is performed),
  For $(2)$, by Lemma~\ref{lem:avss-programtweak},
  the environment can simulate $\Hash$ in a way that is indistinguishable to the adversary
  in the $b=0$ case (the output is chosen at random) and the $b=1$ case (the output is programmed).
  %\ian{As above, there's a negligible advantage here.}
  % CK- I said that in the lemma,we don't need to repeat it again.
  %
  We complete the proof by showing that the adversary has zero advantage for $(3)$.

  Recall that the concrete \AVSS is zero-indexed,
  per Definition~\ref{defn:zero-vss}.
  By Lemma~\ref{lem:avss-simshare}, the environment can embed $\vsschal$ in the first query of the adversary to $\getshare$ in a way that is indistinguishable to $\adv$.
  $\adv$ must query $\getshare$ at least once (i.e., if $\issuedchal=0$ at the end of the experiment, then the experiment returns $0$).
  %\ian{I noted earlier that this is only true if $b=1$; you added a test for $b=1$ to line 12 of the main experiment in Fig 5, but that doesn't change this statement; in fact, now the adversary trivially has perfect advantage by always guessing $b'=0$ and never calling $\getshare$ (though this may depend on the buggy definition of advantage, as earlier). (Separately, I wonder: is it a problem if the adversary never calls $\getauxstring$?)}
  % CK: I changed the definition to always check if a challenge were issued
  % CK: It is not a problem if the adversary did not call getaux; it has negligible chance in winning in that case.
  Hence,
  it is guaranteed that the public recovery set $\{ \vsscommit_j \}_{j \in \set{\total}}$
  that is input to $\avss.\aggregatepub$
  will include a contribution with respect to the discrete logarithm of $\vsschal$.

  Equation~\ref{eq:avss-sim-chal} demonstrates that the commitment to the constant term of $\avsscommit = \{\hat{A}_0, \ldots, \hat{A}_{t-1} \}$
  evaluates to $\vsschal$;
  and so, because the concrete \AVSS is zero-indexed,
  then $\avss.\derivepshare(0, \avsscommit)$ evaluates to $\vsschal$.

\begin{equation}\label{eq:avss-sim-chal}
 \begin{split}
   \hat{A}_0 & = g^{\tweak} \cdot \prod_{j \in \set{\total}} \vsscommit_j[0]   \\
   %
   & = g^\tweak \cdot \vsschal \cdot g^\secretone \cdot g^{\sum_{i \in \set{\simtotal}} \secret_i }  \\
   %
   & = g^{- \secretone - \sum_{i \in \set{\simtotal}} \secret_i } \cdot \vsschal \cdot g^\secretone \cdot g^{\sum_{i \in \set{\simtotal}} \secret_i } \\
   %
   & = \vsschal \\
 \end{split}
\end{equation}

Hence, the adversary gains no advantage in its guess of $b'$,
regardless of whether it is operating in the real ($b=0$)
or simulated ($b=1$) setting.
This concludes the proof.
\end{proof}

\subsection{Uniqueness}

The uniqueness of the concrete \AVSS also follows from that of \FeldVariant, as we now show.

\begin{theorem}
  The concrete \AVSS scheme is unique.
\end{theorem}

\begin{proof}
  We now show that the additional aggregation algorithm in our \AVSS construction gives the adversary no additional power in the uniqueness game than in the plain \FeldVariant setting.

  Recall the VSS uniqueness experiment in Figure~\ref{fg:vss-unique}.
  The adversary is required to output a valid tuple $(\recoveryset_1^*, \recoveryset_2^*, \vsscommit^*)$,
  for which $\vess.\recover(\recoveryset_1^*) \neq \vess.\recover(\recoveryset_2^*)$.
  As above, we know that the adversary cannot win the uniqueness game in the \FeldVariant setting,
  no matter what strategy it employs,
  including using the \AVSS aggregation algorithm,
  since the aggregation algorithm uses no information hidden from the adversary.
  Therefore the adversary cannot win the uniqueness game in the \AVSS setting either.
\end{proof}


